% Introdução: contexto da prática 2
\begin{frame}{Contexto da prática}
  \begin{block}{Problema investigado}
    Determinar como a função de transferência de um filtro ativo influencia sua
    resposta em regime permanente senoidal, confrontando o modelo matemático com
    a simulação e com as medições realizadas em bancada.
  \end{block}

  \begin{columns}[T,onlytextwidth]
    \begin{column}{0.48\textwidth}
      \textbf{Grandezas observadas}
      \begin{itemize}
        \item ganho $|H(j\omega)|=V_o/V_i$;
        \item fase $\phi=360^{\circ}\,\Delta T/T$;
        \item frequência de ressonância;
        \item característica passa-faixa.
      \end{itemize}
    \end{column}
    \begin{column}{0.48\textwidth}
      \textbf{Etapas complementares}
      \begin{itemize}
        \item análise no domínio da frequência;
        \item cálculo simbólico e numérico;
        \item validação no LTspice;
        \item ensaio com protoboard e osciloscópio.
      \end{itemize}
    \end{column}
  \end{columns}
\end{frame}